% Part: many-valued-logic
% Chapter: three-valued-logics
% Section: multiple-designation

\documentclass[../../../include/open-logic-section]{subfiles}

\begin{document}

\olfileid[pt]{mvl}{thr}{mul}

\olsection{Designando não apenas $\True$}

Até aqui, todas as lógicas que vimos tinham como conjunto de valores de verdade
designados $V^+ = \{\True\}$, O.s.\@, algo conta como verdadeiro se, e somente se, seu valor de
verdade é~$\True$. Mas também se poderia contar algo como verdadeiro se simplesmente
não fosse~$\False$. Nesse caso, obteríamos uma lógica estipulando na
matriz, p.ex., que $V^+ = \{\True, \Undef\}$.

\begin{defn}
A \emph{lógica do paradoxo}~$\LogLP$ é definida pela matriz:
\begin{enumerate}
  \item A linguagem proposicional padrão $\Lang L_0$, com
  $\lnot$, $\land$, $\lor$, $\lif$.
  \item O conjunto de valores de verdade $V = \{\True, \Undef, \False\}$.
  \item $\True$ e $\Undef$ são designados, O.s.\@, $V^+ = \{\True, \Undef\}$.
  \item As funções de verdade são as mesmas da lógica forte de Kleene.
\end{enumerate}
\end{defn}

\begin{defn}
A lógica de Halld\'en, a \emph{lógica do absurdo}~$\LogHal$, é definida pela matriz:
\begin{enumerate}
  \item A linguagem proposicional padrão $\Lang L_0$, com
  $\lnot$, $\land$, $\lor$, $\lif$ e um conectivo $1$-ário~$+$.
  \item O conjunto de valores de verdade $V = \{\True, \Undef, \False\}$.
  \item $\True$ e $\Undef$ são designados, O.s.\@, $V^+ = \{\True, \Undef\}$.
  \item As funções de verdade são as mesmas da lógica fraca de Kleene, acrescidas do
  operador ``é sem sentido'':
  \begin{center}
    \begin{tabular}{c|c} 
    $\tf{+}$ & \\ 
    \hline  
    $\True$ & $\False$ \\ 
    $\Undef$ & $\True$ \\
    $\False$ & $\False$ 
    \end{tabular}
  \end{center}
\end{enumerate}
\end{defn}

Ao contrário das lógicas de Kleene, com as quais compartilham tabelas-verdade,
estas \emph{têm} tautologias. 

\begin{prop}\ollabel{prop:LP-taut-CL} As tautologias de $\LogLP$ são
  as mesmas da lógica proposicional clássica.
\end{prop}

\begin{proof}
  Por \olref[syn][sub]{prop:mvl-cl}, se $\Entails[\LogLP] !A$, então
  $\Entails[\LogCL] !A$. Para mostrar a recíproca, mostramos que, se existe
  !!a{valuation}~$\pAssign v\colon \PVar \to \{\False, \True,
  \Undef\}$ tal que $\pValue v(!A)[\LogKs] = \False$, então existe
  !!a{valuation}~$\pAssign {v'}\colon \PVar \to \{\False, \True\}$
  tal que $\pValue {v'}(!A)[\LogCL] = \False$. Isso estabelece o
  resultado para $\LogLP$, pois $\LogKs$ e $\LogLP$ têm as mesmas
  funções de verdade características, e $\False$ é o único valor de verdade
  de $\LogLP$ que não é designado (essa é a única diferença
  entre $\LogLP$ e~$\LogKs$). Assim, se $\Entails/[\LogLP] !A$, para
  alguma !!{valuation}~$\pAssign v$, $\pValue v(!A)[\LogLP] = \pValue
  v[\LogKs](!A) = \False$. Pela afirmação que estamos demonstrando,
  $\pValue{v'}[\LogCL](!A) = \False$, O.s.\@, $\Entails/[\LogCL] !A$.
  
  Para estabelecer a afirmação, primeiro definimos $\pAssign {v'}$ por
  \[
    \pAssign {v'}(p) = 
    \begin{cases}
    \True & \text{se } \pAssign {v}(p) \in \{\True, \Undef\}\\
    \False & \text{caso contrário}
  \end{cases}
  \]
  Mostramos agora, por indução em $!A$, que (a)~se $\pValue v(!A)[\LogKs]
  = \False$, então $\pValue {v'}(!A)[\LogCL] = \False$, e (b)~se
  $\pValue v(!A)[\LogKs] = \True$, então $\pValue {v'}(!A)[\LogCL] =
  \True$.
  \begin{enumerate}
    \item Base da indução: $!A \ident p$. Por
    \olref[syn][val]{defn:pValue}, $\pValue v(!A)[\LogKs]
    = \pAssign v(p) = \pValue {v'}(!A)[\LogCL]$, o que implica tanto (a)
    quanto~(b).
    
    Para o passo de indução, considere os casos:
    \item $!A \ident \lnot !B$. 
    \begin{enumerate}
      \item Suponha $\pValue v(\lnot!B)[\LogKs] = \False$. Pela
      definição de $\tf{\lnot}[\LogKs]$, $\pValue v(!B)[\LogKs]  =
      \True$. Pela hipótese de indução, caso~(b), obtemos $\pValue
      {v'}(!B)[\LogCL]  = \True$; portanto, $\pValue {v'}(\lnot !B)[\LogCL] = \False$.
      \item Suponha $\pValue v(\lnot!B)[\LogKs] = \True$. Pela
      definição de $\tf{\lnot}[\LogKs]$, $\pValue v(!B)[\LogKs]  =
      \False$. Pela hipótese de indução, caso~(a), obtemos $\pValue
      {v'}(!B)[\LogCL]  = \False$; portanto, $\pValue {v'}(\lnot !B)[\LogCL] = \True$.
    \end{enumerate}

    \item $!A \ident (!B \land !C)$.
    \begin{enumerate}
      \item Suponha $\pValue v(!B \land !C)[\LogKs] = \False$. Pela
      definição de $\tf{\land}[\LogKs]$, $\pValue v(!B)[\LogKs]  =
      \False$ ou $\pValue v(!B)[\LogKs]  = \False$. Pela hipótese de
      indução, caso~(a), obtemos $\pValue {v'}(!B)[\LogCL]  = \False$
      ou $\pValue {v'}(!C)[\LogCL]  = \False$; portanto, $\pValue {v'}(!B \land
      !C)[\LogCL] = \False$.
      \item Suponha $\pValue v(!B \land !C)[\LogKs] = \True$. Pela
      definição de $\tf{\land}[\LogKs]$, $\pValue v(!B)[\LogKs]  =
      \True$ e $\pValue v(!B)[\LogKs]  = \True$. Pela hipótese de
      indução, caso~(b), obtemos $\pValue {v'}(!B)[\LogCL]  = \True$
      e $\pValue {v'}(!C)[\LogCL]  = \True$; portanto, $\pValue {v'}(!B \land
      !C)[\LogCL] = \True$.
    \end{enumerate}
  \end{enumerate}

  Os outros dois casos são semelhantes e ficam como exercícios. Alternativamente,
  a prova acima estabelece o resultado para todas as !!{formula}s que contêm
  apenas $\lnot$ e~$\land$. Podemos então recorrer aos fatos de que, tanto em
  $\LogKs$ quanto em $\LogCL$, para qualquer $\pAssign v$, $\pValue v(!B \lor
  !C) = \pValue v(\lnot(\lnot!B \land \lnot !C))$ e $\pValue v(!B \lif
  !C) = \pValue v(\lnot(!B \land \lnot !C))$.
\end{proof}

\begin{prob}
  Complete a prova de \olref[mvl][thr][mul]{prop:LP-taut-CL}, O.s.\@,
  estabeleça (a) e~(b) para os casos em que $!A \ident (!B \lor !C)$
  e $!A \ident (!B \lif !C)$.
\end{prob}

\begin{prob}
Demonstre que toda tautologia clássica é uma tautologia em~$\LogHal$.
\end{prob}

Embora tenham as mesmas tautologias que a lógica clássica, suas
relações de consequência são diferentes. $\LogLP$, por exemplo, é
\emph{paraconsistente}, pois $\lnot p, p \Entails/ q$; assim, o
princípio de explosão $\lnot !A, !A \Entails !B$ não vale em
geral. (Ele vale para alguns casos de $!A$ e $!B$, p.ex., se $!B$~é uma
tautologia.)

\begin{prob}
  Quais das relações a seguir valem (a)~em $\LogLP$
  e (b)~em $\LogHal$? Dê uma tabela-verdade para cada uma.
  \begin{enumerate}
    \item $p, p \lif q \Entails q$
    \item $\lnot q, p \lif q \Entails \lnot p$
    \item $p \lor q, \lnot p \Entails q$
    \item $\lnot p, p \Entails q$
    \item $p \Entails p \lor q$
    \item $p \lif q, q\lif r \Entails p \lif r$
  \end{enumerate}
\end{prob}

O que acontece se designarmos $\Undef$ em $\LogLuk[3]$?

\begin{defn}
  A lógica \emph{R-Mingle trivalorada}~$\LogRM[3]$ é definida pela matriz:
  \begin{enumerate}
    \item A linguagem proposicional padrão $\Lang L_0$, com
    $\lfalse$, $\lnot$, $\land$, $\lor$, $\lif$.
    \item O conjunto de valores de verdade $V = \{\True, \Undef, \False\}$.
    \item $\True$ e $\Undef$ são designados, O.s.\@, $V^+ = \{\True, \Undef\}$.
    \item As funções de verdade são as mesmas da lógica de \L ukasiewicz~$\LogLuk[3]$.
  \end{enumerate}
\end{defn}

\begin{prob}
  Quais das relações a seguir valem em $\LogRM[3]$?
  \begin{enumerate}
    \item $p, p \lif q \Entails q$
    \item $p \lor q, \lnot p \Entails q$
    \item $\lnot p, p \Entails q$
    \item $p \Entails p \lor q$
  \end{enumerate}
\end{prob}

Tabelas-verdade diferentes podem, às vezes, gerar a mesma lógica
(relação de consequência) apenas pela mudança dos valores designados.
P.ex., isso ocorre se, na lógica de G\"odel, tomarmos $V^+ = \{\True,
\Undef\}$ em vez de~$\{\True\}$.

\begin{prop}\ollabel{prop:gl-udes}
  A matriz com $V = \{\False, \Undef,\True\}$, $V^+=\{\True,
  \Undef\}$ e as funções de verdade da lógica $3$-valorada de G\"odel
  define a lógica clássica.
\end{prop}

\begin{proof}
  Exercício. 
\end{proof}

\begin{prob}
Demonstre \olref[mvl][thr][mul]{prop:gl-udes} mostrando que, para a
  lógica~$\Log L$ definida como a lógica de G\"odel, mas com
  $V^+=\{\True,\Undef\}$, se $\Gamma \Entails/[\Log L] !B$, então
  $\Gamma \Entails/[\LogCL] !B$. Use as ideias de
  \olref[mvl][thr][mul]{prop:LP-taut-CL}, mas, em vez de provar
  as propriedades (a) e~(b), mostre que $\pValue v(!A)[\LogGod] = \False$
  se, e somente se, $\pValue {v'}(!A)[\LogCL] = \False$ (e, portanto, que $\pValue
  v(!A)[\LogGod] \in \{\True,\Undef\}$ se, e somente se, $\pValue {v'}(!A)[\LogCL] =
  \True$). Explique por que isso estabelece a proposição.
\end{prob}

\end{document}
